% answers.tex —— 2014 年·广东卷·文科 · 参考答案与详细解析片段
% 由 tools/generate_projects.py 依来源总账自 2014/2014_数学真题_广东文_试题.md + 2014/2014_数学真题_广东文_解析正文.md + 2014/2014_数学真题_广东文_答案要点.md 生成，勿手改（改源 MD 后重跑）


\begin{answer}
\gksolhead{1}{}

\textbf{题目}：

已知集合 $M=\{2,3,4\}$，$N=\{0,2,3,5\}$，则 $M\cap N=$\mbox{（　　）}

（A）$\{0,2\}$；（B）$\{2,3\}$；（C）$\{3,4\}$；（D）$\{3,5\}$。

\Needspace{4\baselineskip}
\textbf{解析}

\Needspace{5\baselineskip}
$M=\{2,3,4\}$，$N=\{0,2,3,5\}$，两集合的公共元素为 $2$ 与 $3$，故
\nopagebreak[4]
\[M\cap N=\{2,3\}.\]

选 B。

\gkanswer{答案：B。$M\cap N=\{2,3,4\}\cap\{0,2,3,5\}=\{2,3\}$。}

\end{answer}
\begin{answer}
\gksolhead{2}{}

\textbf{题目}：

已知复数 $z$ 满足 $(3-4\mathrm i)z=25$，则 $z=$\mbox{（　　）}

（A）$-3-4\mathrm i$；（B）$-3+4\mathrm i$；（C）$3-4\mathrm i$；（D）$3+4\mathrm i$。

\Needspace{4\baselineskip}
\textbf{解析}

\Needspace{5\baselineskip}
两边同乘 $3+4\mathrm i$，利用 $(3-4\mathrm i)(3+4\mathrm i)=9+16=25$，得
\nopagebreak[4]
\[z=\frac{25}{3-4\mathrm i}=\frac{25(3+4\mathrm i)}{25}=3+4\mathrm i.\]

选 D。

\gkanswer{答案：D。$z=\dfrac{25}{3-4\mathrm i}=\dfrac{25(3+4\mathrm i)}{25}=3+4\mathrm i$。}

\end{answer}
\begin{answer}
\gksolhead{3}{}

\textbf{题目}：

已知向量 $\boldsymbol a=(1,2)$，$\boldsymbol b=(3,1)$，则 $\boldsymbol b-\boldsymbol a=$\mbox{（　　）}

（A）$(-2,1)$；（B）$(2,-1)$；（C）$(2,0)$；（D）$(4,3)$。

\Needspace{4\baselineskip}
\textbf{解析}

\Needspace{5\baselineskip}
按分量相减：
\nopagebreak[4]
\[\boldsymbol b-\boldsymbol a=(3,1)-(1,2)=(3-1,1-2)=(2,-1).\]

选 B。

\gkanswer{答案：B。$\boldsymbol b-\boldsymbol a=(3-1,1-2)=(2,-1)$。}

\end{answer}
\begin{answer}
\gksolhead{4}{}

\textbf{题目}：

若变量 $x$，$y$ 满足约束条件 $\begin{cases}x+2y\leqslant 8,\\0\leqslant x\leqslant 4,\\0\leqslant y\leqslant 3,\end{cases}$ 则 $z=2x+y$ 的最大值等于\mbox{（　　）}

（A）$7$；（B）$8$；（C）$10$；（D）$11$。

\Needspace{4\baselineskip}
\textbf{解析}

\Needspace{5\baselineskip}
画出可行域。四条边界 $x+2y=8$，$x=0$，$x=4$，$y=0$，$y=3$ 围成五边形，其顶点为
\nopagebreak[4]
\[(0,0),\quad(4,0),\quad(4,2),\quad(2,3),\quad(0,3).\]
其中 $(4,2)$ 是 $x=4$ 与 $x+2y=8$ 的交点，$(2,3)$ 是 $y=3$ 与 $x+2y=8$ 的交点。

\Needspace{5\baselineskip}
把五个顶点依次代入 $z=2x+y$：
\nopagebreak[4]
\[\begin{aligned}
(0,0)&:\ z=0,\\
(4,0)&:\ z=8,\\
(4,2)&:\ z=2\times4+2=10,\\
(2,3)&:\ z=2\times2+3=7,\\
(0,3)&:\ z=3.
\end{aligned}\]
线性目标函数在可行域上的最大值必在顶点处取得，故 $z_{\max}=10$，在 $(4,2)$ 处取得。

选 C。

\gkanswer{答案：C。可行域顶点 $(0,0)$，$(4,0)$，$(4,2)$，$(2,3)$，$(0,3)$ 处目标值依次为 $0,8,10,7,3$，最大值为 $10$。}

\end{answer}
\begin{answer}
\gksolhead{5}{}

\textbf{题目}：

下列函数为奇函数的是\mbox{（　　）}

（A）$2^x-\dfrac1{2^x}$；（B）$x^3\sin x$；（C）$2\cos x+1$；（D）$x^2+2^x$。

\Needspace{4\baselineskip}
\textbf{解析}

定义域均为 $\mathbf R$，关于原点对称。逐项检验 $f(-x)$ 与 $f(x)$ 的关系：

\Needspace{5\baselineskip}
（A）$f(x)=2^x-\dfrac1{2^x}$，则
\nopagebreak[4]
\[f(-x)=2^{-x}-\frac1{2^{-x}}=2^{-x}-2^{x}=-f(x),\]
为奇函数。

（B）$f(-x)=(-x)^3\sin(-x)=x^3\sin x=f(x)$，为偶函数。

（C）$f(-x)=2\cos(-x)+1=2\cos x+1=f(x)$，为偶函数。

（D）$f(1)=1+2=3$，$f(-1)=1+\dfrac12=\dfrac32$，既不等于 $f(1)$ 也不等于 $-f(1)$，非奇非偶。

故选 A。

\gkanswer{答案：A。$f(x)=2^x-2^{-x}$ 满足 $f(-x)=-f(x)$；（B）（C）为偶函数，（D）非奇非偶。}

\end{answer}
\begin{answer}
\gksolhead{6}{}

\textbf{题目}：

为了解 $1000$ 名学生的学习情况，采用系统抽样的方法，从中抽取容量为 $40$ 的样本，则分段的间隔为\mbox{（　　）}

（A）$50$；（B）$40$；（C）$25$；（D）$20$。

\Needspace{4\baselineskip}
\textbf{解析}

\Needspace{5\baselineskip}
系统抽样把总体均分为 $40$ 段，每段含
\nopagebreak[4]
\[\frac{1000}{40}=25\]
个个体，故分段间隔为 $25$。

选 C。

\gkanswer{答案：C。分段间隔为 $\dfrac{1000}{40}=25$。}

\end{answer}
\begin{answer}
\gksolhead{7}{}

\textbf{题目}：

在 $\triangle ABC$ 中，角 $A$，$B$，$C$ 所对应的边分别为 $a$，$b$，$c$，则“$a\leqslant b$”是“$\sin A\leqslant\sin B$”的\mbox{（　　）}

（A）充分必要条件；（B）充分非必要条件；（C）必要非充分条件；（D）非充分非必要条件。

\Needspace{4\baselineskip}
\textbf{解析}

\Needspace{5\baselineskip}
由正弦定理 $a=2R\sin A$，$b=2R\sin B$（$R>0$），故
\nopagebreak[4]
\[a\leqslant b\iff\sin A\leqslant\sin B.\]
\Needspace{5\baselineskip}
只需再说明两个不等式同解。由和差化积
\nopagebreak[4]
\[\sin A-\sin B=2\cos\frac{A+B}{2}\sin\frac{A-B}{2},\]
\Needspace{5\baselineskip}
而 $A+B=\pi-C$，故
\nopagebreak[4]
\[\cos\frac{A+B}{2}=\cos\left(\frac\pi2-\frac C2\right)=\sin\frac C2>0 .\]
\Needspace{5\baselineskip}
又 $-\dfrac\pi2<\dfrac{A-B}{2}<\dfrac\pi2$，$\sin$ 在该区间上单调递增，于是 $\sin A-\sin B$ 与 $A-B$ 同号，即
\nopagebreak[4]
\[\sin A\leqslant\sin B\iff A\leqslant B .\]
所以“$a\leqslant b$”与“$\sin A\leqslant\sin B$”互为充要条件。

选 A。

\gkanswer{答案：A。由正弦定理 $a\leqslant b\iff\sin A\leqslant\sin B$，且 $\sin A-\sin B=2\cos\dfrac{A+B}{2}\sin\dfrac{A-B}{2}$，$\cos\dfrac{A+B}{2}=\sin\dfrac C2>0$，两者同解。}

\end{answer}
\begin{answer}
\gksolhead{8}{}

\textbf{题目}：

若实数 $k$ 满足 $0<k<5$，则曲线 $\dfrac{x^2}{16}-\dfrac{y^2}{5-k}=1$ 与曲线 $\dfrac{x^2}{16-k}-\dfrac{y^2}{5}=1$ 的\mbox{（　　）}

（A）实半轴长相等；（B）虚半轴长相等；（C）离心率相等；（D）焦距相等。

\Needspace{4\baselineskip}
\textbf{解析}

\Needspace{5\baselineskip}
第一条曲线中 $a_1^2=16$，$b_1^2=5-k$，故
\nopagebreak[4]
\[c_1^2=a_1^2+b_1^2=16+(5-k)=21-k .\]
\Needspace{5\baselineskip}
第二条曲线中 $a_2^2=16-k$，$b_2^2=5$（因 $0<k<5$，两根号内均为正），故
\nopagebreak[4]
\[c_2^2=a_2^2+b_2^2=(16-k)+5=21-k .\]
于是 $c_1^2=c_2^2$，两曲线焦距相等。

再排除其余三项：

实半轴长分别为 $4$ 与 $\sqrt{16-k}$，因 $0<k<5$ 而 $\sqrt{16-k}<4$，不相等；

虚半轴长分别为 $\sqrt{5-k}$ 与 $\sqrt5$，不相等；

离心率分别为 $\dfrac{\sqrt{21-k}}{4}$ 与 $\dfrac{\sqrt{21-k}}{\sqrt{16-k}}$，分母不同，不相等。

选 D。

\gkanswer{答案：D。两曲线 $c^2$ 分别为 $16+(5-k)$ 与 $(16-k)+5$，都等于 $21-k$，故焦距相等。}

\end{answer}
\begin{answer}
\gksolhead{9}{}

\textbf{题目}：

若空间中四条两两不同的直线 $l_1$，$l_2$，$l_3$，$l_4$，满足 $l_1\perp l_2$，$l_2\parallel l_3$，$l_3\perp l_4$，则下列结论一定正确的是\mbox{（　　）}

（A）$l_1\perp l_4$；（B）$l_1\parallel l_4$；（C）$l_1$，$l_4$ 既不垂直也不平行；（D）$l_1$，$l_4$ 的位置关系不确定。

\Needspace{4\baselineskip}
\textbf{解析}

由 $l_1\perp l_2$ 且 $l_2\parallel l_3$ 得 $l_1\perp l_3$。取 $l_3$ 的方向向量为 $(0,1,0)$，则 $l_1$ 的方向向量为 $(1,0,0)$。设 $l_4$ 的方向向量为 $(a,0,c)\neq\boldsymbol 0$，它满足 $l_4\perp l_3$ 且与 $l_1$，$l_3$ 两两不同。

分三种情形：

若 $c=0$，则 $l_4$ 的方向为 $(1,0,0)$，$l_4\parallel l_1$，此时（B）成立；

若 $a=0$，则 $l_4$ 的方向为 $(0,0,1)$，$l_4\perp l_1$，此时（A）成立；

若 $ac\neq0$，取 $l_4$ 过点 $(0,1,1)$，则 $l_4$ 与 $l_1$ 既不平行（方向不成比例）也不垂直（$a\cdot1+0\cdot0+c\cdot0=a\neq0$），也不相交，此时（C）成立。

三种结论都可能出现，故 $l_1$，$l_4$ 的位置关系不确定。

选 D。

\gkanswer{答案：D。$l_1\perp l_2$、$l_2\parallel l_3$ 只给出 $l_1\perp l_3$，而空间中同垂直于 $l_3$ 的两条直线未必平行也未必垂直，故位置关系不确定。}

\end{answer}
\begin{answer}
\gksolhead{10}{}

\textbf{题目}：

对任意复数 $w_1$，$w_2$，定义 $w_1*w_2=w_1\overline{w_2}$，其中 $\overline{w_2}$ 是 $w_2$ 的共轭复数，对任意复数 $z_1$，$z_2$，$z_3$ 有如下四个命题：

① $(z_1+z_2)*z_3=(z_1*z_3)+(z_2*z_3)$；② $z_1*(z_2+z_3)=(z_1*z_2)+(z_1*z_3)$；

③ $(z_1*z_2)*z_3=z_1*(z_2*z_3)$；④ $z_1*z_2=z_2*z_1$。

则真命题的个数是\mbox{（　　）}

（A）$1$；（B）$2$；（C）$3$；（D）$4$。

\Needspace{4\baselineskip}
\textbf{解析}

记 $z_1*z_2=z_1\overline{z_2}$。逐条验证：

① $(z_1+z_2)*z_3=(z_1+z_2)\overline{z_3}=z_1\overline{z_3}+z_2\overline{z_3}=(z_1*z_3)+(z_2*z_3)$，真。

② $z_1*(z_2+z_3)=z_1\overline{(z_2+z_3)}=z_1(\overline{z_2}+\overline{z_3})=(z_1*z_2)+(z_1*z_3)$，真。

\Needspace{5\baselineskip}
③ 左边为
\nopagebreak[4]
\[(z_1*z_2)*z_3=(z_1\overline{z_2})\overline{z_3}=z_1\overline{z_2}\overline{z_3},\]
\Needspace{5\baselineskip}
右边为
\nopagebreak[4]
\[z_1*(z_2*z_3)=z_1\overline{z_2\overline{z_3}}=z_1\overline{z_2}z_3 .\]
取 $z_1=z_2=1$，$z_3=\mathrm i$，则左边 $=-\mathrm i$，右边 $=\mathrm i$，两边不等，假。

④ 左边 $z_1\overline{z_2}$，右边 $z_2\overline{z_1}=\overline{z_1\overline{z_2}}$，二者互为共轭，一般不等。取 $z_1=1$，$z_2=\mathrm i$，则左边 $=-\mathrm i$，右边 $=\mathrm i$，假。

真命题只有 ①②，共 $2$ 个。

选 B。

\gkanswer{答案：B。①②成立；③取 $z_1=z_2=1$，$z_3=\mathrm i$ 时两边为 $-\mathrm i$ 与 $\mathrm i$，④取 $z_1=1$，$z_2=\mathrm i$ 时两边为 $-\mathrm i$ 与 $\mathrm i$，故只有 $2$ 个真命题。}

\end{answer}
\begin{answer}
\gksolhead{11}{}

\textbf{题目}：

曲线 $y=-5\mathrm e^x+3$ 在点 $(0,-2)$ 处的切线方程为\gkblank{}。

\Needspace{4\baselineskip}
\textbf{解析}

\Needspace{5\baselineskip}
$y=-5\mathrm e^x+3$ 的导数为
\nopagebreak[4]
\[y'=-5\mathrm e^x .\]
\Needspace{5\baselineskip}
在 $x=0$ 处 $y'=-5$，即切线斜率为 $-5$，切点为 $(0,-2)$。由点斜式
\nopagebreak[4]
\[y-(-2)=-5(x-0),\]
\Needspace{5\baselineskip}
即
\nopagebreak[4]
\[5x+y+2=0 .\]

\gkanswer{答案：$5x+y+2=0$。$y'=-5\mathrm e^x$，在 $x=0$ 处斜率为 $-5$，切线为 $y+2=-5x$。}

\end{answer}
\begin{answer}
\gksolhead{12}{}

\textbf{题目}：

从字母 $a$，$b$，$c$，$d$，$e$ 中任取两个不同字母，则取到字母 $a$ 的概率为\gkblank{}。

\Needspace{4\baselineskip}
\textbf{解析}

\Needspace{5\baselineskip}
从 $5$ 个字母中任取两个不同字母的取法数为
\nopagebreak[4]
\[\mathrm C_5^2=\frac{5\times4}{2}=10 .\]
\Needspace{5\baselineskip}
含字母 $a$ 的取法：先取定 $a$，再从其余 $4$ 个字母中取 $1$ 个，共
\nopagebreak[4]
\[\mathrm C_4^1=4\]
\Needspace{5\baselineskip}
种。故所求概率为
\nopagebreak[4]
\[P=\frac4{10}=\frac25 .\]

\gkanswer{答案：$\dfrac25$。$\dfrac{\mathrm C_4^1}{\mathrm C_5^2}=\dfrac4{10}=\dfrac25$。}

\end{answer}
\begin{answer}
\gksolhead{13}{}

\textbf{题目}：

等比数列 $\{a_n\}$ 的各项均为正数，且 $a_1a_5=4$，则 $\log_2a_1+\log_2a_2+\log_2a_3+\log_2a_4+\log_2a_5=$\gkblank{}。

\Needspace{4\baselineskip}
\textbf{解析}

$\{a_n\}$ 为各项均为正数的等比数列，故 $a_1a_5=a_3^2=4$，且 $a_3>0$，所以 $a_3=2$。

\Needspace{5\baselineskip}
由等比数列性质 $a_1a_5=a_2a_4=a_3^2$，得
\nopagebreak[4]
\[a_1a_2a_3a_4a_5=a_3^5=2^5=32 .\]
\Needspace{5\baselineskip}
于是
\nopagebreak[4]
\[\log_2a_1+\log_2a_2+\log_2a_3+\log_2a_4+\log_2a_5
=\log_2(a_1a_2a_3a_4a_5)=\log_232=5 .\]

\gkanswer{答案：$5$。由 $a_1a_5=a_3^2=4$ 且各项为正得 $a_3=2$，故 $a_1a_2a_3a_4a_5=a_3^5=32$，所求为 $\log_232=5$。}

\end{answer}
\begin{answer}
\gksolhead{14}{}

\textbf{题目}：

在极坐标系中，曲线 $C_1$ 和 $C_2$ 的方程分别为 $2\rho\cos^2\theta=\sin\theta$ 和 $\rho\cos\theta=1$，以极点为平面直角坐标系的原点，极轴为 $x$ 轴的正半轴，建立平面直角坐标系，则曲线 $C_1$ 和 $C_2$ 的交点的直角坐标为\gkblank{}。

\Needspace{4\baselineskip}
\textbf{解析}

由 $\rho\cos\theta=1$ 得 $x=1$，即 $C_2$ 是直线 $x=1$。

\Needspace{5\baselineskip}
对 $2\rho\cos^2\theta=\sin\theta$ 两边同乘 $\rho$：
\nopagebreak[4]
\[2\rho^2\cos^2\theta=\rho\sin\theta .\]
\Needspace{5\baselineskip}
利用 $x=\rho\cos\theta$，$y=\rho\sin\theta$，得
\nopagebreak[4]
\[2x^2=y,\]
即 $C_1$ 是抛物线 $y=2x^2$。

\Needspace{5\baselineskip}
把 $x=1$ 代入 $y=2x^2$ 得 $y=2$，故交点的直角坐标为
\nopagebreak[4]
\[(1,2).\]

\gkanswer{答案：$(1,2)$。$\rho\cos\theta=1$ 即 $x=1$；$2\rho\cos^2\theta=\sin\theta$ 两边乘 $\rho$ 得 $2x^2=y$，代入 $x=1$ 得 $y=2$。}

\end{answer}
\begin{answer}
\gksolhead{15}{}

\textbf{题目}：

如图，在平行四边形 $ABCD$ 中，点 $E$ 在 $AB$ 上，且 $EB=2AE$，$AC$ 与 $DE$ 交于点 $F$，则 $\dfrac{\triangle CDF\ \text{\text{的周长}}}{\triangle AEF\ \text{\text{的周长}}}=$\gkblank{}。
\nopagebreak[4]
\begin{center}% ============================================================
% f2014_guangdong_arts_q15_parallelogram.tex
% 2014 广东卷（文科）第 15 题：平行四边形 ABCD 中 E 在 AB 上，
% EB=2AE，DE 与 AC 交于 F，求 △CDF 与 △AEF 的周长比。
% 来源 PDF：普通高考/2014/2014广东文.pdf
% 原页：第 1 页（A3 横排）第 2 栏，印面 879，第 15 题题干下方
% 说明：AE:EB=1:2，故 AE=AB/3；图中 A、E、B 依次在底边上，
% D 在 A 上方偏右，C 在 B 上方偏右，DE 近于竖直，AC 为对角线。
% ============================================================
\begin{tikzpicture}[>=Stealth, line width=0.7pt, font=\small,
  x=1cm, y=1cm, line cap=round, line join=round]

  \coordinate (A) at (0,0);
  \coordinate (E) at (2,0);
  \coordinate (B) at (6,0);
  \coordinate (D) at (1.6,3);
  \coordinate (C) at (7.6,3);
  \coordinate (F) at (1.9,0.75);

  \draw (A) -- (B) -- (C) -- (D) -- cycle;
  \draw (D) -- (E);
  \draw (A) -- (C);

  \fill (A) circle (1.1pt) node[below left] {$A$};
  \fill (E) circle (1.1pt) node[below] {$E$};
  \fill (B) circle (1.1pt) node[below right] {$B$};
  \fill (D) circle (1.1pt) node[above left] {$D$};
  \fill (C) circle (1.1pt) node[above right] {$C$};
  \fill (F) circle (1.1pt) node[right] {$F$};
\end{tikzpicture}
\end{center}
\nopagebreak[4]
\Needspace{4\baselineskip}
\textbf{解析}

\Needspace{5\baselineskip}
四边形 $ABCD$ 为平行四边形，故 $AB\parallel DC$，从而 $AE\parallel DC$。又 $A$，$E$，$B$ 共线，$C$，$D$ 在另一条平行线上，$AC$ 与 $DE$ 交于 $F$，故
\nopagebreak[4]
\[\triangle AEF\sim\triangle CDF .\]
\Needspace{5\baselineskip}
相似比为
\nopagebreak[4]
\[\frac{AE}{CD}=\frac{AE}{AB} .\]
\Needspace{5\baselineskip}
由 $EB=2AE$ 得 $AB=AE+EB=3AE$，故
\nopagebreak[4]
\[\frac{AE}{AB}=\frac13 .\]
\Needspace{5\baselineskip}
相似三角形的周长比等于相似比，即
\nopagebreak[4]
\[\frac{\triangle AEF\ \text{\text{的周长}}}{\triangle CDF\ \text{\text{的周长}}}=\frac13,\]
\Needspace{5\baselineskip}
所以
\nopagebreak[4]
\[\frac{\triangle CDF\ \text{\text{的周长}}}{\triangle AEF\ \text{\text{的周长}}}=3 .\]

\gkanswer{答案：$3$。$AE\parallel DC$ 得 $\triangle AEF\sim\triangle CDF$，相似比 $\dfrac{AE}{DC}=\dfrac{AE}{AB}=\dfrac13$，周长比为其倒数 $3$。}

\end{answer}
\begin{answer}
\gksolhead{16}{}

\textbf{题目}：

已知函数 $f(x)=A\sin\left(x+\dfrac\pi3\right)$，$x\in\mathbf R$，且 $f\left(\dfrac{5\pi}{12}\right)=\dfrac{3\sqrt2}{2}$。

（1） 求 $A$ 的值；

（2） 若 $f(\theta)-f(-\theta)=\sqrt3$，$\theta\in\left(0,\dfrac\pi2\right)$，求 $f\left(\dfrac\pi6-\theta\right)$。

\Needspace{4\baselineskip}
\textbf{解析}

\gksubpart{（1）}

\Needspace{5\baselineskip}
由 $f\left(\dfrac{5\pi}{12}\right)=\dfrac{3\sqrt2}{2}$ 得
\nopagebreak[4]
\[A\sin\left(\frac{5\pi}{12}+\frac\pi3\right)=A\sin\frac{3\pi}{4}=\frac{3\sqrt2}{2}.\]
\Needspace{5\baselineskip}
而 $\sin\dfrac{3\pi}{4}=\dfrac{\sqrt2}{2}$，故
\nopagebreak[4]
\[A\cdot\frac{\sqrt2}{2}=\frac{3\sqrt2}{2},\]
解得 $A=3$，即 $f(x)=3\sin\left(x+\dfrac\pi3\right)$。

\gksubpart{（2）}

\Needspace{5\baselineskip}
由 $f(\theta)-f(-\theta)=\sqrt3$ 得
\nopagebreak[4]
\[3\left[\sin\left(\theta+\frac\pi3\right)-\sin\left(-\theta+\frac\pi3\right)\right]=\sqrt3 .\]
\Needspace{5\baselineskip}
用和差化积，或直接展开：
\nopagebreak[4]
\[\begin{aligned}
\sin\left(\theta+\frac\pi3\right)-\sin\left(\frac\pi3-\theta\right)
&=\left(\sin\theta\cos\frac\pi3+\cos\theta\sin\frac\pi3\right)
-\left(\sin\frac\pi3\cos\theta-\cos\frac\pi3\sin\theta\right)\\
&=2\sin\theta\cos\frac\pi3=\sin\theta .
\end{aligned}\]
\Needspace{5\baselineskip}
所以
\nopagebreak[4]
\[3\sin\theta=\sqrt3,\qquad \sin\theta=\frac{\sqrt3}{3}.\]
\Needspace{5\baselineskip}
因 $\theta\in\left(0,\dfrac\pi2\right)$，$\cos\theta>0$，故
\nopagebreak[4]
\[\cos\theta=\sqrt{1-\sin^2\theta}=\sqrt{1-\frac13}=\frac{\sqrt6}{3}.\]
\Needspace{5\baselineskip}
于是
\nopagebreak[4]
\[f\left(\frac\pi6-\theta\right)=3\sin\left(\frac\pi6-\theta+\frac\pi3\right)
=3\sin\left(\frac\pi2-\theta\right)=3\cos\theta=\sqrt6 .\]

\gkanswer{答案（1）：$3$。$\sin\left(\dfrac{5\pi}{12}+\dfrac\pi3\right)=\sin\dfrac{3\pi}4=\dfrac{\sqrt2}{2}$，故 $A\cdot\dfrac{\sqrt2}{2}=\dfrac{3\sqrt2}{2}$。}

\gkanswer{答案（2）：$\sqrt6$。$f(\theta)-f(-\theta)=3\sin\theta=\sqrt3$ 得 $\sin\theta=\dfrac{\sqrt3}{3}$，$\cos\theta=\dfrac{\sqrt6}{3}$，所求为 $3\sin\left(\dfrac\pi2-\theta\right)=3\cos\theta=\sqrt6$。}

\end{answer}
\begin{answer}
\gksolhead{17}{}

\textbf{题目}：

某车间 $20$ 名工人年龄数据如下表：

\begin{center}
\begin{tabular}{ll}
\toprule
年龄（岁） & 工人数（人） \\
\midrule
$19$ & $1$ \\
$28$ & $3$ \\
$29$ & $3$ \\
$30$ & $5$ \\
$31$ & $4$ \\
$32$ & $3$ \\
$40$ & $1$ \\
合计 & $20$ \\
\bottomrule
\end{tabular}
\end{center}

（1） 求这 $20$ 名工人年龄的众数与极差；

（2） 以十位数为茎，个位数为叶，作出这 $20$ 名工人年龄的茎叶图；

（3） 求这 $20$ 名工人年龄的方差。

\Needspace{4\baselineskip}
\textbf{解析}

\gksubpart{（1）}

表中 $30$ 岁出现 $5$ 次，次数最多，故众数为 $30$。

\Needspace{5\baselineskip}
年龄的最大值为 $40$，最小值为 $19$，极差为
\nopagebreak[4]
\[40-19=21 .\]

\gksubpart{（2）}

十位数为茎、个位数为叶，按年龄从小到大排列：

\begin{center}
\begin{tabular}{ll}
\toprule
茎 & 叶 \\
\midrule
$1$ & $9$ \\
$2$ & $8\ 8\ 8\ 9\ 9\ 9$ \\
$3$ & $0\ 0\ 0\ 0\ 0\ 1\ 1\ 1\ 1\ 2\ 2\ 2$ \\
$4$ & $0$ \\
\bottomrule
\end{tabular}
\end{center}

\gksubpart{（3）}

\Needspace{5\baselineskip}
各年龄的总和为
\nopagebreak[4]
\[19\times1+28\times3+29\times3+30\times5+31\times4+32\times3+40\times1
=19+84+87+150+124+96+40=600 .\]
\Needspace{5\baselineskip}
平均年龄为
\nopagebreak[4]
\[\bar x=\frac{600}{20}=30 .\]
\Needspace{5\baselineskip}
各年龄与平均年龄之差的平方和为
\nopagebreak[4]
\[\begin{aligned}
&(19-30)^2\times1+(28-30)^2\times3+(29-30)^2\times3+(30-30)^2\times5\\
&\quad +(31-30)^2\times4+(32-30)^2\times3+(40-30)^2\times1\\
={}&121+12+3+0+4+12+100=252 .
\end{aligned}\]
\Needspace{5\baselineskip}
故方差为
\nopagebreak[4]
\[s^2=\frac{252}{20}=12.6 .\]

\gkanswer{答案（1）：众数为 $30$，极差为 $21$。$30$ 岁出现 $5$ 次最多；最大 $40$、最小 $19$。}

\gkanswer{答案（2）：茎 $1$ 叶 $9$；茎 $2$ 叶 $8,8,8,9,9,9$；茎 $3$ 叶 $0,0,0,0,0,1,1,1,1,2,2,2$；茎 $4$ 叶 $0$。}

\gkanswer{答案（3）：$12.6$。平均年龄 $\dfrac{600}{20}=30$，偏差平方和 $252$，方差 $\dfrac{252}{20}=12.6$。}

\end{answer}
\begin{answer}
\gksolhead{18}{}

\textbf{题目}：

如图 1，四边形 $ABCD$ 为矩形，$PD\perp$ 平面 $ABCD$，$AB=1$，$BC=PC=2$，作如图 2 折叠，折痕 $EF\parallel DC$。其中点 $E$，$F$ 分别在线段 $PD$，$PC$ 上，沿 $EF$ 折叠后点 $P$ 叠在线段 $AD$ 上的点记为 $M$，并且 $MF\perp CF$。
\nopagebreak[4]
\begin{center}% ============================================================
% f2014_guangdong_arts_q18_fold.tex
% 2014 广东卷（文科）第 18 题：矩形 ABCD 中 PD⊥平面 ABCD，
% AB=1，BC=PC=2，沿 EF（EF∥DC，E∈PD，F∈PC）折叠，P 落在
% 线段 AD 上的 M 处，且 MF⊥CF。左为折叠前（图 1），右为折叠后（图 2）。
% 来源 PDF：普通高考/2014/2014广东文.pdf
% 原页：第 1 页（A3 横排）第 3 栏下部，印面 879，第 18 题题干下方
% 说明：图 2 依原页保留折叠前 P 点及 P—E、P—F 两段（示意折痕两侧），
% 并以虚线表示被遮挡的棱 MD、DC、DE、EC；三棱锥 M-CDE 的 ME、MC 为实线。
% ============================================================
\begin{tikzpicture}[>=Stealth, line width=0.7pt, font=\small,
  x=1cm, y=1cm, line cap=round, line join=round]

  % ---------------- 图 1：折叠前 ----------------
  \coordinate (A1) at (0,4.6);
  \coordinate (B1) at (3.4,4.6);
  \coordinate (C1) at (3.72,1.0);
  \coordinate (D1) at (0.3,1.0);
  \coordinate (P1) at (-1.0,-1.0);

  \draw (A1) -- (B1) -- (C1) -- (D1) -- cycle;
  \draw (P1) -- (D1);
  \draw (P1) -- (C1);

  \fill (A1) circle (1.1pt) node[above left] {$A$};
  \fill (B1) circle (1.1pt) node[above right] {$B$};
  \fill (C1) circle (1.1pt) node[right] {$C$};
  \fill (D1) circle (1.1pt) node[above left] {$D$};
  \fill (P1) circle (1.1pt) node[below left] {$P$};
  \node at (1.7,-1.35) {图 1};

  % ---------------- 图 2：折叠后 ----------------
  \coordinate (A2) at (6.2,4.6);
  \coordinate (B2) at (9.6,4.6);
  \coordinate (C2) at (9.92,1.0);
  \coordinate (D2) at (6.5,1.0);
  \coordinate (P2) at (5.2,-1.0);
  \coordinate (M2) at (6.3164,3.2032);
  \coordinate (E2) at (6.175,0.5);
  \coordinate (F2) at (8.74,0.5);

  % 矩形边框（DC 被折起部分遮挡）
  \draw (A2) -- (B2) -- (C2);
  \draw (A2) -- (M2);
  \draw[dashed] (M2) -- (D2) -- (C2);

  % 被遮挡的棱
  \draw[dashed] (E2) -- (D2);
  \draw[dashed] (E2) -- (C2);

  % 折痕与折叠后的三角形 MEF
  \draw (E2) -- (F2);
  \draw (F2) -- (C2);
  \draw (M2) -- (E2);
  \draw (M2) -- (F2);
  \draw (M2) -- (C2);

  % 折叠前 P 的位置（示意）
  \draw[line width=0.55pt] (P2) -- (E2);
  \draw[line width=0.55pt] (P2) -- (F2);

  \fill (A2) circle (1.1pt) node[above left] {$A$};
  \fill (B2) circle (1.1pt) node[above right] {$B$};
  \fill (C2) circle (1.1pt) node[right] {$C$};
  \fill (D2) circle (1.1pt) node[above right, inner sep=2pt] {$D$};
  \fill (M2) circle (1.1pt) node[left] {$M$};
  \fill (E2) circle (1.1pt) node[left] {$E$};
  \fill (F2) circle (1.1pt) node[below right] {$F$};
  \fill (P2) circle (1.1pt) node[below left] {$P$};
  \node at (7.9,-1.35) {图 2};
\end{tikzpicture}
\end{center}
\nopagebreak[4]
（1） 证明：$CF\perp$ 平面 $MDF$；

（2） 求三棱锥 $M-CDE$ 的体积。

\Needspace{4\baselineskip}
\textbf{解析}

\gksubpart{（1）}

\Needspace{5\baselineskip}
因为 $PD\perp$ 平面 $ABCD$，而 $AD\subset$ 平面 $ABCD$，所以
\nopagebreak[4]
\[PD\perp AD .\]
又 $ABCD$ 为矩形，所以 $AD\perp DC$。

\Needspace{5\baselineskip}
$M$ 在线段 $AD$ 上，故 $MD$ 与 $AD$ 共线，于是
\nopagebreak[4]
\[MD\perp PD,\qquad MD\perp DC .\]
\Needspace{5\baselineskip}
$PD$ 与 $DC$ 是平面 $PDC$ 内两条相交直线，所以
\nopagebreak[4]
\[MD\perp\text{\text{平面} }PDC .\]
\Needspace{5\baselineskip}
而 $CF\subset$ 平面 $PDC$，故
\nopagebreak[4]
\[MD\perp CF .\]
\Needspace{5\baselineskip}
又已知 $MF\perp CF$，且 $MD\cap MF=M$，$MD$，$MF\subset$ 平面 $MDF$，所以
\nopagebreak[4]
\[CF\perp\text{\text{平面} }MDF .\]

\gksubpart{（2）}

\Needspace{5\baselineskip}
由（1）$CF\perp$ 平面 $MDF$，而 $DF\subset$ 平面 $MDF$，所以
\nopagebreak[4]
\[CF\perp DF .\]

\Needspace{5\baselineskip}
在 $\mathrm{Rt}\triangle PDC$ 中，$PD\perp DC$，$DC=AB=1$，$PC=2$，故
\nopagebreak[4]
\[PD=\sqrt{PC^2-DC^2}=\sqrt{4-1}=\sqrt3 .\]
\Needspace{5\baselineskip}
设
\nopagebreak[4]
\[PE=\sqrt3\,t,\qquad PF=2t\qquad(0<t<1),\]
由 $EF\parallel DC$ 得 $\dfrac{PE}{PD}=\dfrac{PF}{PC}=t$，且 $EF=t\cdot DC=t$。

\Needspace{5\baselineskip}
于是
\nopagebreak[4]
\[CF=PC-PF=2-2t,\qquad DE=PD-PE=\sqrt3(1-t).\]

\Needspace{5\baselineskip}
在 $\triangle DCF$ 中，$\angle DCF=\angle DCP$，而
\nopagebreak[4]
\[\cos\angle DCP=\frac{DC}{PC}=\frac12 .\]
\Needspace{5\baselineskip}
由余弦定理
\nopagebreak[4]
\[\begin{aligned}
DF^2&=DC^2+CF^2-2\cdot DC\cdot CF\cos\angle DCF\\
&=1+(2-2t)^2-2\times1\times(2-2t)\times\frac12\\
&=1+4(1-t)^2-2(1-t).
\end{aligned}\]
\Needspace{5\baselineskip}
由 $CF\perp DF$ 得 $CF^2+DF^2=DC^2=1$，即
\nopagebreak[4]
\[4(1-t)^2+1+4(1-t)^2-2(1-t)=1,\]
\Needspace{5\baselineskip}
化简得
\nopagebreak[4]
\[8(1-t)^2-2(1-t)=0,\]
\Needspace{5\baselineskip}
即 $2(1-t)\left[4(1-t)-1\right]=0$。因 $0<t<1$，故
\nopagebreak[4]
\[1-t=\frac14,\qquad t=\frac34 .\]
\Needspace{5\baselineskip}
于是
\nopagebreak[4]
\[PF=\frac32,\qquad CF=\frac12,\qquad DF=\sqrt{1-\frac14}=\frac{\sqrt3}{2},\qquad DE=\frac{\sqrt3}{4}.\]

\Needspace{5\baselineskip}
折叠保持长度，故 $MF=PF=\dfrac32$。由 $MD\perp$ 平面 $PDC$ 得 $MD\perp DF$，在 $\mathrm{Rt}\triangle MDF$ 中
\nopagebreak[4]
\[MD=\sqrt{MF^2-DF^2}=\sqrt{\frac94-\frac34}=\sqrt{\frac32}=\frac{\sqrt6}{2}.\]

\Needspace{5\baselineskip}
因为 $C$，$D$，$E$ 都在平面 $PDC$ 内，所以平面 $CDE$ 就是平面 $PDC$，点 $M$ 到该平面的距离即为 $MD=\dfrac{\sqrt6}{2}$。又 $DC\perp DE$，故
\nopagebreak[4]
\[S_{\triangle CDE}=\frac12\cdot DC\cdot DE=\frac12\times1\times\frac{\sqrt3}{4}=\frac{\sqrt3}{8}.\]
\Needspace{5\baselineskip}
所以
\nopagebreak[4]
\[V_{M-CDE}=\frac13S_{\triangle CDE}\cdot MD
=\frac13\times\frac{\sqrt3}{8}\times\frac{\sqrt6}{2}
=\frac{\sqrt2}{16}.\]

\gkanswer{答案（1）：见解析版}

\gkanswer{答案（2）：$\dfrac{\sqrt2}{16}$。由 $CF\perp$ 平面 $MDF$ 得 $CF\perp DF$，解得 $PF=\dfrac32$、$CF=\dfrac12$、$DF=\dfrac{\sqrt3}{2}$、$MD=\dfrac{\sqrt6}{2}$；$S_{\triangle CDE}=\dfrac{\sqrt3}{8}$，$V=\dfrac13\cdot\dfrac{\sqrt3}{8}\cdot\dfrac{\sqrt6}{2}=\dfrac{\sqrt2}{16}$。}

\end{answer}
\begin{answer}
\gksolhead{19}{}

\textbf{题目}：

设各项均为正数的数列 $\{a_n\}$ 的前 $n$ 项和为 $S_n$，且 $S_n$ 满足 $S_n^2-(n^2+n-3)S_n-3(n^2+n)=0$，$n\in\mathbf N^*$。

（1） 求 $a_1$ 的值；

（2） 求数列 $\{a_n\}$ 的通项公式；

（3） 证明：对一切正整数 $n$，有 $\dfrac1{a_1(a_1+1)}+\dfrac1{a_2(a_2+1)}+\cdots+\dfrac1{a_n(a_n+1)}<\dfrac13$。

\Needspace{4\baselineskip}
\textbf{解析}

\gksubpart{（1）}

\Needspace{5\baselineskip}
在 $S_n^2-(n^2+n-3)S_n-3(n^2+n)=0$ 中令 $n=1$。此时 $S_1=a_1$，$n^2+n=2$，得
\nopagebreak[4]
\[a_1^2-(1+1-3)a_1-3\times2=0,\]
\Needspace{5\baselineskip}
即
\nopagebreak[4]
\[a_1^2+a_1-6=0,\]
\Needspace{5\baselineskip}
分解得 $(a_1+3)(a_1-2)=0$。因各项均为正数，$a_1>0$，故
\nopagebreak[4]
\[a_1=2 .\]

\gksubpart{（2）}

\Needspace{5\baselineskip}
把方程左边分解。令 $u=n^2+n$，则方程为
\nopagebreak[4]
\[S_n^2-(u-3)S_n-3u=0,\]
\Needspace{5\baselineskip}
而
\nopagebreak[4]
\[(S_n-u)(S_n+3)=S_n^2+(3-u)S_n-3u,\]
\Needspace{5\baselineskip}
正好与左边相同，故
\nopagebreak[4]
\[(S_n-n^2-n)(S_n+3)=0 .\]
\Needspace{5\baselineskip}
因为 $a_n>0$，所以 $S_n>0$，从而 $S_n+3\neq0$，只能
\nopagebreak[4]
\[S_n=n^2+n .\]
\Needspace{5\baselineskip}
当 $n\geqslant2$ 时，
\nopagebreak[4]
\[\begin{aligned}
a_n&=S_n-S_{n-1}\\
&=(n^2+n)-\left[(n-1)^2+(n-1)\right]\\
&=n^2+n-(n^2-n)=2n .
\end{aligned}\]
\Needspace{5\baselineskip}
又 $a_1=2=2\times1$，符合上式，故
\nopagebreak[4]
\[a_n=2n\qquad(n\in\mathbf N^*).\]

\gksubpart{（3）}

\Needspace{5\baselineskip}
由（2）得 $a_k=2k$，故
\nopagebreak[4]
\[\frac1{a_k(a_k+1)}=\frac1{2k(2k+1)} .\]
当 $n=1$ 时，左边 $=\dfrac16<\dfrac13$，结论成立。

\Needspace{5\baselineskip}
以下设 $n\geqslant2$。对 $k\geqslant2$，
\nopagebreak[4]
\[2k(2k+1)=4k^2+2k>4k^2-1=(2k-1)(2k+1),\]
\Needspace{5\baselineskip}
故
\nopagebreak[4]
\[\frac1{2k(2k+1)}<\frac1{(2k-1)(2k+1)}
=\frac12\left(\frac1{2k-1}-\frac1{2k+1}\right).\]
\Needspace{5\baselineskip}
于是
\nopagebreak[4]
\[\begin{aligned}
\sum_{k=1}^{n}\frac1{a_k(a_k+1)}
&\leqslant\frac16+\frac12\sum_{k=2}^{n}\left(\frac1{2k-1}-\frac1{2k+1}\right)\\
&<\frac16+\frac12\left(\frac13-\frac1{2n+1}\right)
<\frac16+\frac16=\frac13 .
\end{aligned}\]
故对一切正整数 $n$，所证不等式成立。

\gkanswer{答案（1）：$2$。令 $n=1$ 得 $a_1^2+a_1-6=0$，取正根 $a_1=2$。}

\gkanswer{答案（2）：$a_n=2n$。方程分解为 $(S_n-n^2-n)(S_n+3)=0$，故 $S_n=n^2+n$，相减得 $a_n=2n$。}

\gkanswer{答案（3）：见解析版}

\end{answer}
\begin{answer}
\gksolhead{20}{}

\textbf{题目}：

已知椭圆 $C:\dfrac{x^2}{a^2}+\dfrac{y^2}{b^2}=1$（$a>b>0$）的一个焦点为 $(\sqrt5,0)$，离心率为 $\dfrac{\sqrt5}{3}$。

（1） 求椭圆 $C$ 的标准方程；

（2） 若动点 $P(x_0,y_0)$ 为椭圆 $C$ 外一点，且点 $P$ 到椭圆 $C$ 的两条切线相互垂直，求点 $P$ 的轨迹方程。

\Needspace{4\baselineskip}
\textbf{解析}

\gksubpart{（1）}

\Needspace{5\baselineskip}
设焦距参数为 $c$。由焦点 $(\sqrt5,0)$ 在 $x$ 轴上得 $c=\sqrt5$。由离心率
\nopagebreak[4]
\[e=\frac ca=\frac{\sqrt5}{3}\]
\Needspace{5\baselineskip}
得
\nopagebreak[4]
\[a=\frac{c}{e}=\frac{\sqrt5}{\sqrt5/3}=3,\]
\Needspace{5\baselineskip}
故
\nopagebreak[4]
\[b^2=a^2-c^2=9-5=4 .\]
\Needspace{5\baselineskip}
椭圆的标准方程为
\nopagebreak[4]
\[\frac{x^2}{9}+\frac{y^2}{4}=1 .\]

\gksubpart{（2）}

\Needspace{5\baselineskip}
设 $P(x_0,y_0)$ 为椭圆外一点，先设过 $P$ 的切线斜率存在，记为 $k$，则切线为
\nopagebreak[4]
\[y=kx+m,\qquad m=y_0-kx_0 .\]
\Needspace{5\baselineskip}
代入椭圆方程 $4x^2+9y^2=36$：
\nopagebreak[4]
\[4x^2+9(kx+m)^2=36,\]
\Needspace{5\baselineskip}
即
\nopagebreak[4]
\[(4+9k^2)x^2+18km\,x+9m^2-36=0 .\]
\Needspace{5\baselineskip}
切线与椭圆只有一个公共点，故判别式为零：
\nopagebreak[4]
\[(18km)^2-4(4+9k^2)(9m^2-36)=0 .\]
\Needspace{5\baselineskip}
两边除以 $36$ 得
\nopagebreak[4]
\[9k^2m^2-(4+9k^2)(m^2-4)=0,\]
\Needspace{5\baselineskip}
展开并消去 $9k^2m^2$：
\nopagebreak[4]
\[36k^2+16-4m^2=0,\]
\Needspace{5\baselineskip}
即
\nopagebreak[4]
\[m^2=9k^2+4 .\]
\Needspace{5\baselineskip}
把 $m=y_0-kx_0$ 代入：
\nopagebreak[4]
\[(y_0-kx_0)^2=9k^2+4,\]
\Needspace{5\baselineskip}
整理成关于 $k$ 的二次方程
\nopagebreak[4]
\[(9-x_0^2)k^2+2x_0y_0\,k+(4-y_0^2)=0 .\]
\Needspace{5\baselineskip}
两条切线的斜率 $k_1$，$k_2$ 是该方程的两根（$9-x_0^2\neq0$ 时），由韦达定理
\nopagebreak[4]
\[k_1k_2=\frac{4-y_0^2}{9-x_0^2} .\]
\Needspace{5\baselineskip}
两切线互相垂直，故 $k_1k_2=-1$，即
\nopagebreak[4]
\[\frac{4-y_0^2}{9-x_0^2}=-1,\]
\Needspace{5\baselineskip}
解得
\nopagebreak[4]
\[x_0^2+y_0^2=13 .\]

再补上切线斜率不存在的情形：此时 $x_0=\pm3$，竖直线 $x=x_0$ 与椭圆相切，另一条切线必须水平，斜率为 $0$，即 $y=y_0$ 与椭圆相切，故 $y_0=\pm2$。此时仍有 $x_0^2+y_0^2=9+4=13$，与上式一致。

\Needspace{5\baselineskip}
反过来，若 $x_0^2+y_0^2=13$，则
\nopagebreak[4]
\[\frac{x_0^2}{9}+\frac{y_0^2}{4}\geqslant\frac{x_0^2+y_0^2}{9}=\frac{13}{9}>1,\]
点 $P$ 在椭圆外，上述推导可逆，两条切线互相垂直。

\Needspace{5\baselineskip}
故点 $P$ 的轨迹方程为
\nopagebreak[4]
\[x^2+y^2=13 .\]

\gkanswer{答案（1）：$\dfrac{x^2}{9}+\dfrac{y^2}{4}=1$。$c=\sqrt5$，由 $\dfrac ca=\dfrac{\sqrt5}{3}$ 得 $a=3$，$b^2=9-5=4$。}

\gkanswer{答案（2）：$x^2+y^2=13$。设过 $P$ 的切线 $y=kx+m$，相切条件化为 $(9-x_0^2)k^2+2x_0y_0k+(4-y_0^2)=0$，由 $k_1k_2=-1$ 得 $x_0^2+y_0^2=13$（竖直切线情形同样满足）。}

\end{answer}
\begin{answer}
\gksolhead{21}{}

\textbf{题目}：

已知函数 $f(x)=\dfrac13x^3+x^2+ax+1$（$a\in\mathbf R$）。

（1） 求函数 $f(x)$ 的单调区间；

（2） 当 $a<0$ 时，试讨论是否存在 $x_0\in\left(0,\dfrac12\right)\cup\left(\dfrac12,1\right)$，使得 $f(x_0)=f\left(\dfrac12\right)$。

\Needspace{4\baselineskip}
\textbf{解析}

\gksubpart{（1）}

\Needspace{5\baselineskip}
函数的定义域为 $\mathbf R$，导数为
\nopagebreak[4]
\[f'(x)=x^2+2x+a=(x+1)^2+a-1 .\]

当 $a\geqslant1$ 时，$f'(x)\geqslant0$ 恒成立（仅 $a=1$ 且 $x=-1$ 时取等号），故 $f(x)$ 在 $(-\infty,+\infty)$ 上单调递增。

\Needspace{5\baselineskip}
当 $a<1$ 时，$f'(x)=0$ 有两个不同实根
\nopagebreak[4]
\[x_1=-1-\sqrt{1-a},\qquad x_2=-1+\sqrt{1-a},\]
且 $x_1<x_2$。$f'(x)$ 是开口向上的二次函数，故

在 $(-\infty,x_1)$ 与 $(x_2,+\infty)$ 上 $f'(x)>0$，$f(x)$ 单调递增；

在 $(x_1,x_2)$ 上 $f'(x)<0$，$f(x)$ 单调递减。

\Needspace{5\baselineskip}
综上：$a\geqslant1$ 时单调递增区间为 $(-\infty,+\infty)$；$a<1$ 时增区间为
\nopagebreak[4]
\[\left(-\infty,-1-\sqrt{1-a}\right),\qquad\left(-1+\sqrt{1-a},+\infty\right),\]
\Needspace{5\baselineskip}
减区间为
\nopagebreak[4]
\[\left(-1-\sqrt{1-a},-1+\sqrt{1-a}\right).\]

\gksubpart{（2）}

\Needspace{5\baselineskip}
计算 $f(x)-f\left(\dfrac12\right)$。由
\nopagebreak[4]
\[f\left(\frac12\right)=\frac13\cdot\frac18+\frac14+\frac a2+1=\frac{31}{24}+\frac a2,\]
\Needspace{5\baselineskip}
得
\nopagebreak[4]
\[\begin{aligned}
f(x)-f\left(\frac12\right)
&=\frac13\left(x^3-\frac18\right)+\left(x^2-\frac14\right)+a\left(x-\frac12\right)\\
&=\left(x-\frac12\right)\left[\frac13\left(x^2+\frac x2+\frac14\right)+\left(x+\frac12\right)+a\right]\\
&=\left(x-\frac12\right)\left(\frac13x^2+\frac76x+\frac7{12}+a\right).
\end{aligned}\]
\Needspace{5\baselineskip}
记
\nopagebreak[4]
\[\varphi(x)=\frac13x^2+\frac76x+\frac7{12}+a .\]
\Needspace{5\baselineskip}
因 $x_0\neq\dfrac12$，条件 $f(x_0)=f\left(\dfrac12\right)$ 等价于
\nopagebreak[4]
\[\varphi(x_0)=0,\qquad x_0\in\left(0,\frac12\right)\cup\left(\frac12,1\right).\]

\Needspace{5\baselineskip}
$\varphi(x)$ 是开口向上的二次函数，其对称轴为
\nopagebreak[4]
\[x=-\frac{7/6}{2/3}=-\frac74<0,\]
\Needspace{5\baselineskip}
故 $\varphi(x)$ 在 $(0,1)$ 上单调递增。于是 $\varphi$ 在 $(0,1)$ 内至多一个零点，且

\nopagebreak[4]
\[\varphi(0)=\frac7{12}+a,\qquad
\varphi\left(\frac12\right)=\frac1{12}+\frac7{12}+\frac7{12}+a=\frac54+a,\qquad
\varphi(1)=\frac13+\frac76+\frac7{12}+a=\frac{25}{12}+a .\]

按 $a<0$ 讨论：

\Needspace{5\baselineskip}
若 $a\in\left(-\dfrac54,-\dfrac7{12}\right)$，则
\nopagebreak[4]
\[\varphi(0)=\frac7{12}+a<0,\qquad \varphi\left(\frac12\right)=\frac54+a>0,\]
由单调性，存在唯一 $x_0\in\left(0,\dfrac12\right)$ 使 $\varphi(x_0)=0$，符合要求。

\Needspace{5\baselineskip}
若 $a\in\left(-\dfrac{25}{12},-\dfrac54\right)$，则
\nopagebreak[4]
\[\varphi\left(\frac12\right)=\frac54+a<0,\qquad \varphi(1)=\frac{25}{12}+a>0,\]
同理存在唯一 $x_0\in\left(\dfrac12,1\right)$ 使 $\varphi(x_0)=0$，符合要求。

若 $a=-\dfrac54$，则 $\varphi\left(\dfrac12\right)=0$，零点恰为 $x_0=\dfrac12$，不在所给区间内；而 $\varphi$ 单调递增，别处无零点，故不存在。

若 $a\leqslant-\dfrac{25}{12}$，则 $\varphi(1)=\dfrac{25}{12}+a\leqslant0$，由单调性 $\varphi(x)<0$ 对一切 $x\in(0,1)$ 成立（$a=-\dfrac{25}{12}$ 时零点为 $x=1$），故不存在。

若 $a\geqslant-\dfrac7{12}$，则 $\varphi(0)=\dfrac7{12}+a\geqslant0$，由单调性 $\varphi(x)>0$ 对一切 $x\in(0,1)$ 成立（$a=-\dfrac7{12}$ 时零点为 $x=0$），故不存在。

\Needspace{5\baselineskip}
综上，当
\nopagebreak[4]
\[-\frac{25}{12}<a<-\frac7{12}\quad\text{\text{且}}\quad a\neq-\frac54\]
时存在满足条件的 $x_0$；当 $a\leqslant-\dfrac{25}{12}$、$a=-\dfrac54$ 或 $-\dfrac7{12}\leqslant a<0$ 时不存在。

\gkanswer{答案（1）：当 $a\geqslant1$ 时，单调递增区间为 $(-\infty,+\infty)$；当 $a<1$ 时，增区间为 $(-\infty,-1-\sqrt{1-a})$ 与 $(-1+\sqrt{1-a},+\infty)$，减区间为 $(-1-\sqrt{1-a},-1+\sqrt{1-a})$。}

\gkanswer{答案（2）：当 $-\dfrac{25}{12}<a<-\dfrac{7}{12}$ 且 $a\neq-\dfrac54$ 时存在；当 $a\leqslant-\dfrac{25}{12}$、$a=-\dfrac54$ 或 $-\dfrac{7}{12}\leqslant a<0$ 时不存在。依据是 $f(x)-f\left(\dfrac12\right)=\left(x-\dfrac12\right)\varphi(x)$，$\varphi(x)=\dfrac13x^2+\dfrac76x+\dfrac7{12}+a$ 在 $(0,1)$ 上递增，须 $\varphi$ 在 $(0,1)$ 内有零点且零点不为 $\dfrac12$。}

\end{answer}